%% F09 --- Vectors in the plane and in space
%%
%% Every number in this program that the reader cannot do in their head comes
%% from code/f09_vectors.py and is pulled in with \val{}. (1,2)+(3,4) = (4,6),
%% 3(1,2) = (3,6) and |(3,4)| = 5 are written inline, because they are the
%% arithmetic the reader is here to do.
%%
%% THE THESIS: a vector is a LIST OF NUMBERS you may ALSO draw as an arrow.
%% The drawing is a convenience of two and three dimensions rather than part
%% of the definition, which is why the arithmetic survives into 768 dimensions
%% unchanged and the picture does not.
%%
%% THE PAYOFF IS A BRIDGE and it is an identity, not a table:
%%   |a - b|^2 = |a|^2 + |b|^2 - 2 (a . b)
%% which is F08's a.b = |a||b|cos(theta) rearranged. Two consequences, both
%% swept in the script:
%%   on the UNIT SPHERE, |a - b|^2 = 2 - 2 cos(theta), so distance is a
%%     strictly decreasing function of the cosine and the two measures CANNOT
%%     rank differently -- checked on 13026 comparisons;
%%   off it they can, and the frames work one concrete triple where the
%%     nearest neighbour by distance is not the most similar by cosine.
%%
%% WHAT THIS PROGRAM DELIBERATELY LEAVES ALONE, checked against
%% tools/programs.json rather than remembered:
%%   vector spaces, span, independence, basis, dimension          -> P04
%%   inner products in general, projection, L1 against L2, the
%%     disagreement worked out, what normalising costs you, and
%%     near-orthogonality in high dimension                       -> P05
%%   broadcasting, reshape, axes                                  -> P07
%% F09 owns the arithmetic and the identity that says when the two measures
%% cannot disagree. It hands over everything that needs a vector space.
%%
%% Twin: programs/pl/F09-vectors.tex. Frame for frame, macro for macro,
%% number for number; tools/parity.py compares them.

\program{Vectors in the plane and in space}
\label{prog:F09}
\index{vector!as a list of numbers}
\index{vector!length and unit vector}

\begin{outcomes}
\outcome{1--6}{say what a vector is without drawing one, add and subtract two
  of them, and multiply one by a number}
\outcome{7--13}{compute a length in two and in three dimensions, and turn a
  vector into one of length $1$}
\outcome{14--18}{say what the dot product computes and connect it to the angle
  Program~\ref{prog:F08} found}
\outcome{19--27}{write the distance between two vectors in terms of their
  lengths and their dot product, and say what that identity settles about
  distance and cosine similarity}
\outcome{28--32}{say which parts of the two-dimensional picture survive into
  seven hundred dimensions and which are gone}
\end{outcomes}

\begin{quiz}
\item What is $(1, 2) + (3, 4)$? \teachesat{1--2}
  \answerto{$(4, 6)$. Addition is done one component at a time and nothing
  else.}
\item What is $3(1, 2)$? \teachesat{3--4}
  \answerto{$(3, 6)$. Multiplying by a number multiplies every component.}
\item What is the length of $(3, 4)$? \teachesat{7--8}
  \answerto{$5$, from $\sqrt{3^{2} + 4^{2}}$, which is Pythagoras.}
\item What is the length of $(2, 3, 6)$? \teachesat{8--9}
  \answerto{$\val{f09.len3d}$, from $\sqrt{2^{2} + 3^{2} + 6^{2}}$. Pythagoras
  applied twice.}
\item How do you turn a vector into one of length $1$? \teachesat{11--12}
  \answerto{Divide every component by its length. Every vector except the zero
  vector can be done this way.}
\item Which vector has no direction? \teachesat{12--13}
  \answerto{The zero vector. Normalising it divides by zero, and nothing else
  in this program has an exception.}
\item What does $a \cdot b$ compute, in words? \teachesatone{14}
  \answerto{Multiply the two vectors component by component and add up the
  results.}
\item Two vectors have lengths $3$ and $4$. What can you say about the length
  of their sum? \teachesat{19--20}
  \answerto{It is at most $7$, and it is $7$ only when the two point the same
  way.}
\item Write the squared distance $\lVert a - b \rVert^{2}$ using only lengths
  and a dot product. \teachesatone{21}
  \answerto{$\lVert a \rVert^{2} + \lVert b \rVert^{2} - 2\,(a \cdot b)$.}
\item Two vectors both have length $1$. If one is nearer to a query than the
  other, is it also the more similar by cosine? \teachesat{22--23}
  \answerto{Yes, always. On the unit sphere the distance is
  $2 - 2\cos\theta$, which falls as the cosine rises, so the two orderings are
  the same one.}
\item And if the lengths are not $1$? \teachesat{24--25}
  \answerto{Then they can disagree, and the program works a case where they
  do.}
\item How many components does an embedding of dimension
  $\val{f09.dim}$ have, and how many of them can you draw?
  \teachesat{29--30}
  \answerto{$\val{f09.dim}$ and $\val{f09.drawable}$.}
\item Which of these survives into $\val{f09.dim}$ dimensions: the arrow, the
  arithmetic, or both? \teachesat{28--30}
  \answerto{The arithmetic. Every operation in this program is a loop over the
  components and never once needed the picture.}
\item Name one thing you would be wrong to assume about
  $\val{f09.dim}$-dimensional vectors on the strength of a two-dimensional
  drawing. \teachesatone{31}
  \answerto{Any answer that names a claim about how vectors are typically
  arranged rather than about the arithmetic. The arithmetic transfers; the
  intuition about what is near what does not, and
  Program~\ref{prog:P05} measures the difference.}
\end{quiz}

% ============================================================
\section{A list you may also draw}
\label{sec:F09-list}
\index{vector!as a list of numbers}

\begin{fr}
A \textbf{vector} is a list of numbers.

That is the definition, and it is worth taking it exactly as literally as it
is written, because everything else in this program is arithmetic on lists and
none of it needs anything more.

When the list has two numbers in it you may \emph{also} draw it: put a dot at
those coordinates and an arrow from the origin to the dot. When it has three
you may still draw it, with more difficulty. When it has
$\val{f09.dim}$ you may not, and the arithmetic does not notice.

Two vectors: $a = (1, 2)$ and $b = (3, 4)$. What do you think $a + b$ is?
\nextframe
\end{fr}

\begin{fr}
\begin{ansblock}
$(4, 6)$
\end{ansblock}

One component at a time: $1 + 3$ and $2 + 4$. There is nothing else it could
sensibly mean, and the picture agrees \dash{} lay $b$'s arrow at the tip of
$a$'s and you arrive at $(4, 6)$.

Note which way round that went. The arithmetic came first and the picture
agreed with it. That order matters more than it looks, and section 5 says why.

Subtraction is the same, one component at a time. What is $b - a$?
\nextframe
\end{fr}

\begin{fr}
\ans{$(2, 2)$}
And that one has a reading worth keeping: $b - a$ is the arrow that takes you
\emph{from} $a$ \emph{to} $b$. Every distance in this program is going to be
the length of a difference, so this is the first half of that idea.

The third operation is multiplying by a plain number \dash{} a
\textbf{scalar}, which is the word for a number when there are vectors about.
What is $3(1, 2)$?
\nextframe
\end{fr}

\begin{fr}
\ans{$(3, 6)$}
Every component multiplied. The arrow points the same way and is three times
as long, so a scalar changes the size and not the direction.

Two special cases are worth naming now because both come back:

\begin{itemize}
\item a \textbf{negative} scalar reverses the arrow as well as scaling it, so
  $-1$ times a vector points the opposite way;
\item \textbf{zero} times any vector is the zero vector, all of whose
  components are $0$, and which is the one vector with no direction at all.
\end{itemize}

Those three operations \dash{} add, subtract, scale \dash{} are the whole of
the algebra. What is $2(1, 2) - (3, 4)$?
\nextframe
\end{fr}

\begin{fr}
\ans{$(-1, 0)$}
From $(2, 4) - (3, 4)$. Scale first, then subtract; the conventions are the
ones Program~\ref{prog:F02} set for numbers and nothing new is being asked of
you.

\begin{trapbox}
\textbf{Adding a plain number to a vector is not defined}, and this is worth
saying because an array library will do it anyway. In the mathematics of this
program $v + 1$ means nothing: $v$ is a list and $1$ is not, and there is no
component-by-component pairing to make.

An array library will interpret it as adding $1$ to \emph{every} component,
which is a reasonable convention and is not what the mathematics says. It has
a name and a set of rules, and Program~\ref{prog:P07} states them. Until then,
the useful habit is to notice when you have written a line whose meaning comes
from the library rather than from the mathematics \dash{} because that is the
line whose behaviour changes when the shapes do.
\end{trapbox}

\mermaidfig{f09-list-or-arrow}{The definition is the list. The arrow is a
convenience that runs out.}{a vector as a list, and when the arrow stops}
\end{fr}

\begin{fr}
One more piece of vocabulary, and it is one the book will use constantly.

The numbers in the list are its \textbf{components}, and the number of them is
how many the vector has \dash{} two for a point in the plane,
$\val{f09.dim}$ for the embeddings you will meet. When two vectors have
different numbers of components, none of the three operations above is
defined, because there is nothing to pair the leftover components with.

That is the entire content of a shape error, and it is why one is an error
rather than a wrong answer.
\end{fr}

% ============================================================
\section{Length}
\label{sec:F09-length}
\index{vector!length and unit vector}

\begin{fr}
A vector drawn as an arrow has a length, and the length is Pythagoras.

For $(3, 4)$: go $3$ across and $4$ up, and the arrow is the hypotenuse of a
right-angled triangle with those two sides. So its length is
$\sqrt{3^{2} + 4^{2}}$.

Work that out.
\nextframe
\end{fr}

\begin{fr}
\ans{$5$}
And the notation for it is $\lVert a \rVert$, with double bars, read
\enquote{the length of $a$} or \enquote{the norm of $a$}. You have already
seen it: Program~\ref{prog:F08} put it in the denominator of the
cosine-similarity formula and named it the length without saying how to
compute one, and this is how.

Now three components. What is the length of $(2, 3, 6)$?
\dotline
\nextframe
\end{fr}

\begin{fr}
\ans{$\val{f09.len3d}$}
From $\sqrt{2^{2} + 3^{2} + 6^{2}} = \sqrt{4 + 9 + 36} = \sqrt{49}$.

The reason it is the same formula with one more square in it is worth seeing
once, because it is what makes the generalisation honest rather than assumed.
Take the arrow to $(2, 3, 6)$. Its shadow on the floor is the arrow to
$(2, 3)$, whose length is $\sqrt{2^{2} + 3^{2}}$ by the previous frame. That
shadow and the height $6$ are the two sides of a second right-angled triangle,
and its hypotenuse is the arrow itself:

\[ \lVert (2, 3, 6) \rVert = \sqrt{\left(\sqrt{2^{2} + 3^{2}}\right)^{2} + 6^{2}}
   = \sqrt{2^{2} + 3^{2} + 6^{2}} \]

\textbf{Pythagoras applied twice collapses into one sum of squares}, and the
same argument applied a fourth time, and a fifth, gives the same shape again.
So the length of a list of any number of components is the square root of the
sum of their squares, and nothing in the derivation ever needed to see more
than one triangle at a time.

\mermaidfig{f09-length-twice}{One more component is one more square in the
sum, and the argument is the same triangle again.}{length in two dimensions
and in three}
\end{fr}

\begin{fr}
Try it at a size you cannot draw. A vector with $\val{f09.dim}$ components,
every one of them $1$. What is its length?
\nextframe
\end{fr}

\begin{fr}
\ans{$\val{f09.len.ones}$}
Because the sum of squares is $\val{f09.dim}$ ones, so the length is
$\sqrt{\val{f09.dim}}$.

That is the first calculation in this book that could not have been drawn, and
it went through without incident. Hold on to how unremarkable it was.

Now the operation that uses the length. A \textbf{unit vector} is one whose
length is $1$. How would you make one out of $(3, 4)$?
\nextframe
\end{fr}

\begin{fr}
\begin{ansblock}
Divide by its length: $\left(\frac{3}{5}, \frac{4}{5}\right)$.
\end{ansblock}

Dividing every component by $\lVert a \rVert$ scales the arrow by
$\frac{1}{\lVert a \rVert}$, which by the previous section changes the size and
not the direction \dash{} so the result points the same way and has length
exactly $1$. This is called \textbf{normalising}, and $\hat{a}$ is the usual
name for the result.

Every vector can be normalised except one. Which, and why?
\nextframe
\end{fr}

\begin{fr}
\begin{ansblock}
The zero vector: dividing by its length divides by zero.
\end{ansblock}

And that is not an accident of the arithmetic. The zero vector genuinely has
no direction \dash{} there is no way it points \dash{} so there is nothing for
a unit vector to be.

\begin{warning}
A library will tell you about this in its own way, and the way depends on the
library. Plain Python raises \code{ZeroDivisionError} and stops. An array
library is more likely to produce components that are all \code{nan} and carry
on, with a warning you may not be reading, and then every downstream
comparison against those components is \code{False} \dash{} including the ones
that ask whether it is equal to itself.

The invariant is what to remember: \textbf{the zero vector cannot be
normalised}, whatever your tools do about it. A pipeline that normalises
embeddings needs an answer for the empty document, and \enquote{whatever the
library does} is not one.
\end{warning}
\end{fr}

% ============================================================
\section{The dot product}
\label{sec:F09-dot}
\index{vector!dot product}

\begin{fr}
The fourth operation is the one that turns two vectors into a single number.

The \textbf{dot product} of $a$ and $b$ is written $a \cdot b$ and computed by
multiplying the components in pairs and adding up:
\[ a \cdot b = \sum_i a_i b_i \]
which is Program~\ref{prog:F04}'s sigma with a product inside it, and no more
than that.

Compute $(1, 2) \cdot (3, 4)$.
\nextframe
\end{fr}

\begin{fr}
\ans{$11$}
From $1 \times 3 + 2 \times 4 = 3 + 8$, which is the sum
Program~\ref{prog:F04} works when it asks whether a sigma distributes over a
product \dash{} the same computation, now with a name. Note the shape of the
answer: two vectors in, one plain number out. Nothing about it is a vector, which is why
it is sometimes called the \emph{scalar} product.

One special case is worth taking now, because it will save you writing.
What is $a \cdot a$?
\nextframe
\end{fr}

\begin{fr}
\begin{ansblock}
The sum of the squares of the components, which is $\lVert a \rVert^{2}$.
\end{ansblock}

So $\lVert a \rVert = \sqrt{a \cdot a}$, and the length is not a separate idea
after all \dash{} it is the dot product of a vector with itself, with a square
root on the outside. Every formula in the rest of this program that has a
length in it could have been written with a dot product instead, and the next
section takes advantage of that.

Now the connection you already have. Program~\ref{prog:F08} showed that
cosine similarity is the cosine of the angle between two vectors:
\[ \cos\theta = \frac{a \cdot b}{\lVert a \rVert\, \lVert b \rVert} \]
Multiply both sides by the two lengths and read the result.
\nextframe
\end{fr}

\begin{fr}
\begin{ansblock}
\[ a \cdot b = \lVert a \rVert\, \lVert b \rVert \cos\theta \]
\end{ansblock}

\textbf{The dot product is the two lengths multiplied together and then scaled
by how aligned the two are.} That is the whole meaning of the operation, and
it is why one number can carry information about two vectors at once: it
carries their sizes and their alignment multiplied into each other.

Three cases read straight off it, and all three are worth being able to say
without computing anything:

\begin{center}
\begin{tabular}{lcc}
\toprule
The two point & $\cos\theta$ & $a \cdot b$ \\
\midrule
the same way & $1$ & $\lVert a \rVert\, \lVert b \rVert$ \\
at right angles & $0$ & $0$ \\
opposite ways & $-1$ & $-\lVert a \rVert\, \lVert b \rVert$ \\
\bottomrule
\end{tabular}
\end{center}

The middle row is the one that earns its keep: \textbf{a dot product of zero
means the two are at right angles}, whatever their lengths, and that is how
the word \emph{orthogonal} is usually being used when you meet it.

Two vectors have lengths $2$ and $5$ and a dot product of $0$. What is the
angle between them?
\nextframe
\end{fr}

\begin{fr}
\ans{A right angle}
Because $\cos\theta = \frac{0}{2 \times 5} = 0$, and the lengths never
mattered. A zero dot product says the same thing whatever the two vectors are
scaled to.
\end{fr}

% ============================================================
\section{Distance, and when it agrees with similarity}
\label{sec:F09-distance}
\index{vector!distance against similarity}

\begin{fr}
Two vectors, two questions you might ask about them.

\textbf{How far apart are they?} Subtract, and take the length of the
difference: $\lVert a - b \rVert$. This is the \textbf{Euclidean distance},
and it is what a query in a vector database may be measuring.

\textbf{How alike are their directions?} Divide the dot product by both
lengths: $\frac{a \cdot b}{\lVert a \rVert \lVert b \rVert}$. This is cosine
similarity, and it is what the same query may be measuring instead.

Before anything else, one thing about the first. Vectors $a$ and $b$ have
lengths $3$ and $4$. What is the largest $\lVert a + b \rVert$ could be?
\nextframe
\end{fr}

\begin{fr}
\ans{$\val{f09.tri.sum}$}
And it gets there only when the two point the same way, so that laying one
arrow at the tip of the other walks in a straight line. Any other angle and
the walk turns a corner, and a corner is a short cut:
\[ \lVert a + b \rVert \le \lVert a \rVert + \lVert b \rVert \]

This is the \textbf{triangle inequality}, and its content is that the direct
route is never longer than going via a third point. At right angles \dash{}
$(3, 0)$ and $(0, 4)$ \dash{} the sum has length $\val{f09.tri.actual}$
against the $\val{f09.tri.sum}$ you would get by adding the lengths, so a
right angle costs two units of the seven.

\begin{trapbox}
\textbf{Lengths do not add, and the mistake is to carry the arithmetic of
positive numbers across without checking that it survives.} Two lengths of $3$
and $4$ laid end to end on a line do measure $7$, which is where the habit
comes from; a vector carries a direction as well as a size, and as soon as the
directions differ the two arrows no longer lie end to end. Equality is now the
exception rather than the rule \dash{} it happens only when the two point the
same way.

That is the same shape as Program~\ref{prog:F02}'s expansion trap and
Program~\ref{prog:F06}'s inequality trap: a rule that is true in one setting,
generalised past its hypotheses because nothing announced that the setting had
changed. It is worth expecting the next time an operation gets a new kind of
argument.

Checked in \code{code/f09\_vectors.py} on every pair from a grid of
twenty-five by twenty-five, with equality reproducing to
$\val{f09.tri.parallel.err}$ for parallel pairs.
\end{trapbox}
\end{fr}

\begin{fr}
Now put the two questions together, which is what the rest of this section is
about.

Start from $\lVert a - b \rVert^{2}$ and use the fact that a squared length is
a dot product with itself. Expanding gives
\[ \lVert a - b \rVert^{2} = (a - b) \cdot (a - b)
   = \lVert a \rVert^{2} + \lVert b \rVert^{2} - 2\,(a \cdot b) \]
because the middle terms are $-a \cdot b$ twice.

That identity is checked over sixty-eight thousand pairs in the script, with a
worst disagreement of $\val{f09.law.err}$. Now suppose both vectors have been
normalised, so both lengths are $1$. What does it become?
\nextframe
\end{fr}

\begin{fr}
\begin{ansblock}
\[ \lVert a - b \rVert^{2} = 2 - 2\cos\theta \]
\end{ansblock}

Because both squared lengths are $1$, and $a \cdot b$ is $\cos\theta$ once the
lengths are gone \dash{} which is Program~\ref{prog:F08}'s result with the
denominator equal to $1$. Checked to $\val{f09.unit.err}$.

Read what that says. On the right there is nothing but the cosine, and it
comes with a minus sign. So as the cosine goes up the distance goes down, and
as the cosine goes down the distance goes up, with no exceptions and no
regions where the behaviour changes.

You have a query and two candidates, all three normalised. Candidate A is
nearer to the query than candidate B. Which has the higher cosine similarity?
\dotline
\nextframe
\end{fr}

\begin{fr}
\ans{A, necessarily}
\textbf{On the unit sphere, distance and cosine similarity are one question
with two answers.} Not usually, not approximately: they rank every pair of
candidates identically, because one is a strictly decreasing function of the
other and a strictly decreasing function cannot reorder anything. The script
checks it on $\val{f09.rank.checked}$ comparisons and finds no exception,
which is what you would expect from an identity.

This is worth having because it settles an argument you will meet. If your
vectors are normalised, \enquote{should we use cosine or Euclidean distance?}
has no answer, because there is no difference to have an opinion about.

\mermaidfig{f09-two-questions}{Two questions that become one as soon as both
vectors have length one.}{distance and similarity on the unit sphere}
\end{fr}

\begin{fr}
And now take the normalising away, which is where the trouble is.

A query points along $(1, 0)$. Two candidates:
\[ A = (\num{0.30},\ \num{0.06}) \qquad B = (\num{0.95},\ \num{0.45}) \]

$A$ points almost exactly the way the query does and is short. $B$ points less
the same way but is about the query's own size.

Which is the query's nearest neighbour, and which is its most similar? Compute
both, for both, before reading on.
\dotline
\nextframe
\end{fr}

\begin{fr}
\begin{ansblock}
\begin{center}
\begin{tabular}{lcc}
\toprule
& cosine & distance \\
\midrule
$A$ & $\val{f09.dis.cos.a}$ & $\val{f09.dis.dist.a}$ \\
$B$ & $\val{f09.dis.cos.b}$ & $\val{f09.dis.dist.b}$ \\
\bottomrule
\end{tabular}
\end{center}
\end{ansblock}

\textbf{$A$ is the more similar and $B$ is the nearer.} The two measures name
different winners on the same three vectors, and neither is wrong: they were
asked different questions and they answered them.

\begin{trapbox}
\textbf{\enquote{Nearest} and \enquote{most similar} are the same ranking only
on the unit sphere, and the previous frame is exactly why.} Take the lengths
away and $\lVert a - b \rVert^{2} = 2 - 2\cos\theta$ becomes
$\lVert a \rVert^{2} + \lVert b \rVert^{2} - 2\lVert a \rVert \lVert b \rVert
\cos\theta$, which has the two lengths in it in their own right. The cosine no
longer determines the distance, so it can no longer order it.

What makes this expensive rather than interesting is that both queries run,
both return results that look plausible, and the difference shows up as
retrieval that is slightly wrong in a way nobody can reproduce from the
formula. \textbf{The question to ask of any similarity index is not which
measure it uses but whether the vectors going into it were normalised}, because
the first question only has consequences when the answer to the second is no.

Program~\ref{prog:P05} takes the comparison further \dash{} which measure to
prefer when they do disagree, and what normalising costs you.
\end{trapbox}
\end{fr}

\begin{fr}
One last check on the same three vectors, because it makes the fix concrete
rather than a slogan.

Normalise $A$ and $B$, leave the query alone \dash{} it already has length $1$
\dash{} and measure the distances again. What has to happen?
\nextframe
\end{fr}

\begin{fr}
\begin{ansblock}
$A$ must become the nearer, because the cosines have not changed and now they
determine the distances.
\end{ansblock}

And they do: $\val{f09.dis.ndist.a}$ against $\val{f09.dis.ndist.b}$, where
before normalising it was $\val{f09.dis.dist.a}$ against
$\val{f09.dis.dist.b}$ the other way round.

Nothing about the \emph{directions} changed \dash{} normalising cannot change
a direction, by section 2 \dash{} so the cosines are the same two numbers they
were. What changed is that the distance is now a function of them alone.

\begin{aibox}
This is the whole of why an embedding pipeline normalises before it indexes,
and it is a design decision rather than a formality. Normalised vectors make
the index's choice of metric irrelevant, so the same store answers a cosine
query and a distance query identically and cannot be misconfigured into
disagreeing with itself.

What it costs is the lengths, and the lengths may carry information \dash{}
how much text went into the embedding, how confident the model was. If they
do, normalising discards it before the index ever sees it, which is a decision
worth taking deliberately rather than by copying a tutorial.
\end{aibox}
\end{fr}

% ============================================================
\section{Where the picture runs out}
\label{sec:F09-highdim}
\index{vector!in high dimension}

\begin{fr}
Count what this program has actually used.

Addition, subtraction and scaling are one loop over the components. Length is
a sum of squares and a square root. The dot product is a sum of products.
Distance is a subtraction and a length. Every one of them is a pass over a
list, and the number of components appears in exactly one place: how many
times round the loop.

So which of these operations gets harder when the vector has
$\val{f09.dim}$ components rather than $2$?
\nextframe
\end{fr}

\begin{fr}
\begin{ansblock}
None of them. Each takes more steps and none takes a different kind of step.
\end{ansblock}

The arithmetic never needed the picture, which is why it survives. You drew
arrows in this program because they made the two-dimensional case easy to
believe, not because the definitions used them \dash{} and every definition
was stated as a rule about components before any arrow was drawn.

Now the other side. Which of these operations can you still \emph{draw} at
$\val{f09.dim}$ components?
\nextframe
\end{fr}

\begin{fr}
\begin{ansblock}
None of them either. $\val{f09.drawable}$ components can be drawn and
$\val{f09.undrawn}$ cannot.
\end{ansblock}

\textbf{The arithmetic transfers and the picture does not}, and this is the
sentence the program exists to leave you with.

An embedding is a point in a space with too many dimensions to draw. Every
formula in this program applies to it exactly, with no adjustment. What does
not apply is your sense of where things are: that came from drawings of two
dimensions, and it was never derived from anything.
\end{fr}

\begin{fr}
It is worth being precise about which claims are in which class, because the
distinction is easy to state and hard to hold on to.

\begin{center}
\begin{tabularx}{\linewidth}{@{}lX@{}}
\toprule
Transfers exactly & Does not transfer \\
\midrule
every identity in this program & how many directions are far apart \\
the triangle inequality & how much of the space is near the middle \\
distance and cosine agreeing on the unit sphere & what a
  \enquote{typical} pair of vectors looks like \\
\bottomrule
\end{tabularx}
\end{center}

The left column is proved and holds at any number of components. The right
column is about how vectors are \emph{arranged}, which is a question about
distributions rather than about the operations, and none of it has been
established here.

\begin{warning}
\textbf{The right-hand column is where high dimension surprises people, and
this program deliberately does not tell you the answers.}
Program~\ref{prog:P05} measures them, and the reason for the delay is that
measuring them properly needs the vocabulary Program~\ref{prog:P04} sets up.

What you can take from here is the discipline rather than a result: when you
find yourself reasoning about $\val{f09.dim}$-dimensional embeddings by
picturing two arrows on a page, notice which column your claim is in. If it is
an identity, the picture is a legitimate aid. If it is about what is near what,
the picture is evidence about a plane and you are talking about somewhere else.
\end{warning}
\end{fr}

\begin{fr}
One closing observation, and it is about how this program was built rather
than about vectors.

Every definition here was given as an operation on a list, and the arrow came
afterwards as a check that the operation was sensible. It would have been
easier to do it the other way round \dash{} define everything by pictures and
then say the formulas compute it \dash{} and the result would have been a
program that stopped working at three components.

\textbf{When a definition and a picture agree, it is worth knowing which one
is load-bearing}, because that is what tells you where the definition still
applies once the picture is gone.
\end{fr}

% ============================================================
\begin{summarybox}
\sumitem{1--4}{\result{\textbf{A vector is a list of numbers.} That is the whole
  definition; the arrow is a convenience available at two and three components
  and gone after that. Addition, subtraction and multiplication by a scalar
  are done one component at a time, and none of them needs the picture}.}
\sumitem{5--6}{\result{Adding a plain number to a vector is \textbf{not defined} in
  the mathematics, whatever an array library does with it
  (Program~\ref{prog:P07} has the rules). And two vectors with different
  numbers of components support none of the three operations, which is the
  entire content of a shape error}.}
\sumitem{7--9}{\result{$\lVert a \rVert$ is the square root of the sum of the squares,
  and it is Pythagoras applied once per component: the length in three
  dimensions is built from the length in two, and the same argument runs at
  any number. $\lVert (3,4) \rVert = 5$ and
  $\lVert (2,3,6) \rVert = \val{f09.len3d}$}.}
\sumitem{10--13}{\result{\textbf{Normalising divides every component by the length}
  and gives a vector of length $1$ pointing the same way, because scaling
  changes size and not direction. Every vector can be normalised except the
  zero vector, which has no direction for a unit vector to have \dash{} and
  what your tools do about that case is a property of your tools, not of the
  mathematics}.}
\sumitem{14--16}{\result{$a \cdot b = \sum_i a_i b_i$ takes two vectors and returns
  one number. $a \cdot a = \lVert a \rVert^{2}$, so length is not a separate
  idea \dash{} it is a dot product with a square root on it}.}
\sumitem{17--18}{\result{\textbf{$a \cdot b = \lVert a \rVert \lVert b \rVert
  \cos\theta$}, which is Program~\ref{prog:F08}'s result rearranged: the dot
  product carries both sizes and the alignment multiplied together. A dot
  product of \emph{zero} means a right angle whatever the lengths, and that is
  what \emph{orthogonal} means}.}
\sumitem{19--20}{\result{\textbf{Lengths do not add.}
  $\lVert a + b \rVert \le \lVert a \rVert + \lVert b \rVert$, with equality
  only when the two point the same way; at right angles $3$ and $4$ give
  $\val{f09.tri.actual}$ rather than $\val{f09.tri.sum}$}.}
\sumitem{21--22}{\result{\textbf{$\lVert a - b \rVert^{2} = \lVert a \rVert^{2} +
  \lVert b \rVert^{2} - 2\,(a \cdot b)$}, checked to $\val{f09.law.err}$, and
  on the unit sphere it collapses to $2 - 2\cos\theta$, checked to
  $\val{f09.unit.err}$}.}
\sumitem{23}{\result{So \textbf{on the unit sphere distance and cosine similarity rank
  identically} \dash{} one is a strictly decreasing function of the other and
  cannot reorder anything. Confirmed on $\val{f09.rank.checked}$ comparisons.
  If your vectors are normalised, \enquote{cosine or Euclidean?} has no
  answer}.}
\sumitem{24--25}{\result{\textbf{Off the sphere they can name different winners.} On
  one query and two candidates, $A$ is the more similar
  ($\val{f09.dis.cos.a}$ against $\val{f09.dis.cos.b}$) and $B$ is the nearer
  ($\val{f09.dis.dist.b}$ against $\val{f09.dis.dist.a}$). Both queries run,
  both return plausible results, and the disagreement is invisible in the
  output. The question to ask of an index is not which measure it uses but
  whether what went into it was normalised}.}
\sumitem{26--27}{\result{Normalising the same two candidates puts $A$ nearer
  ($\val{f09.dis.ndist.a}$ against $\val{f09.dis.ndist.b}$) without changing
  either direction \dash{} which is why a pipeline normalises before it
  indexes, and what it pays is the lengths}.}
\sumitem{28--30}{\result{Every operation here is one pass over the components, so
  \textbf{the arithmetic transfers to $\val{f09.dim}$ dimensions unchanged and
  the picture does not}: $\val{f09.drawable}$ components can be drawn and
  $\val{f09.undrawn}$ cannot}.}
\sumitem{31--32}{\result{Identities transfer exactly. Claims about \emph{where things
  are} \dash{} what is near what, what a typical pair looks like \dash{} are
  about how vectors are arranged rather than about the operations, and none of
  them is established by a drawing of two arrows.
  Program~\ref{prog:P05} measures them. When a definition and a picture agree,
  know which one is load-bearing}.}
\end{summarybox}

\canyou

% ============================================================
\begin{testexercises}
\item Compute $(2, -1) + (0, 5)$ and $(2, -1) - (0, 5)$.
  \answerto{$(2, 4)$ and $(2, -6)$. One component at a time, both times.}
\item Compute $-2(3, -1, 4)$.
  \answerto{$(-6, 2, -8)$. Every component multiplied, and the sign reverses
  the direction.}
\item Give the length of $(0, -5)$ and of $(1, 1, 1, 1)$.
  \answerto{$5$ and $2$. The second is $\sqrt{4}$, and negative components
  make no difference because they are squared.}
\item Normalise $(0, -5)$.
  \answerto{$(0, -1)$. Divide both components by $5$; the direction is
  unchanged, so it still points downwards.}
\item A vector has $\val{f09.dim}$ components, every one of them $1$. Give its
  length, and say what changes if every component is $-1$ instead.
  \answerto{$\val{f09.len.ones}$ both times. Squaring removes the sign, so the
  length is the same and only the direction has reversed.}
\item Compute $(2, 3) \cdot (3, -2)$ and say what it tells you.
  \answerto{$0$, so the two are at right angles. It says nothing about their
  lengths.}
\item Two vectors have lengths $6$ and $\num{0.5}$ and point in exactly the
  same direction. Give their dot product.
  \answerto{$3$, from $\lVert a \rVert \lVert b \rVert \cos\theta$ with
  $\cos\theta = 1$.}
\item Two vectors have lengths $2$ and $3$. Give the largest and the smallest
  their dot product could be.
  \answerto{$6$ and $-6$, at $\cos\theta = 1$ and $\cos\theta = -1$ \dash{}
  the same way and opposite ways.}
\item Vectors of lengths $5$ and $12$ are at right angles. Give the length of
  their sum.
  \answerto{$13$, from $\lVert a + b \rVert^{2} = 25 + 144 - 0$. Not $17$,
  which is what adding the lengths would give.}
\item Two unit vectors are at right angles. Use
  $\lVert a - b \rVert^{2} = 2 - 2\cos\theta$ to give the distance between
  them.
  \answerto{$\sqrt{2}$, because $\cos\theta = 0$.}
\item A store holds normalised embeddings and is switched from a cosine index
  to a Euclidean one. What changes in the results?
  \answerto{Nothing. On the unit sphere the two rank identically, so every
  query returns the same neighbours in the same order.}
\item The same store holds unnormalised embeddings and the same switch is
  made. What changes?
  \answerto{The ranking can change, and nothing in the output says so. Long
  vectors are penalised by distance and ignored by cosine.}
\item Which of these is a claim the two-dimensional picture is evidence for at
  $\val{f09.dim}$ components: that the triangle inequality holds; that two
  vectors picked at random are usually far apart in angle?
  \answerto{The first. It is an identity and holds at any number of
  components. The second is about how vectors are arranged and a drawing of a
  plane says nothing about it.}
\item Explain in one sentence why every operation in this program costs more
  at $\val{f09.dim}$ components but none of them gets harder.
  \answerto{Each is a single pass over the components, so the number of
  components sets how many times round the loop and never what happens
  inside it.}
\end{testexercises}

% ============================================================
\begin{furtherproblems}
\item Show that $\lVert ka \rVert = \lvert k \rvert \lVert a \rVert$ for any
  scalar $k$, and say why the modulus signs are needed.
  \answerto{Each component becomes $k a_i$, so each square becomes
  $k^{2} a_i^{2}$ and $k^{2}$ comes out of the root as $\lvert k \rvert$. A
  length cannot be negative, and $k$ can be.}
\item Two vectors satisfy $\lVert a + b \rVert = \lVert a \rVert +
  \lVert b \rVert$. What do you know about them?
  \answerto{One is a non-negative multiple of the other \dash{} they point the
  same way, or one of them is zero. Equality in the triangle inequality
  happens nowhere else.}
\item Show that $\lVert a + b \rVert^{2} + \lVert a - b \rVert^{2} =
  2\lVert a \rVert^{2} + 2\lVert b \rVert^{2}$.
  \answerto{Expand both squares with the identity of frame $21$: the
  $-2(a \cdot b)$ and the $+2(a \cdot b)$ cancel and the four squared lengths
  remain.}
\item Two unit vectors are at an angle of $60$ degrees. Give the distance
  between them, and notice what is special about the answer.
  \answerto{$1$, from $\sqrt{2 - 2 \times \num{0.5}}$. The two unit vectors
  and their difference form an equilateral triangle.}
\item Suppose every vector in a store is multiplied by $10$. What happens to
  every cosine similarity, and what happens to every distance?
  \answerto{Every cosine is unchanged, because scaling both vectors leaves the
  angle alone. Every distance is multiplied by $10$. The rankings by each
  measure are unchanged, and the two may still disagree with each other
  exactly as before.}
\item A colleague reports that normalising the embeddings \enquote{fixed} a
  retrieval bug. Give the most likely mechanism, and say what evidence would
  confirm it.
  \answerto{The index was ranking by distance while the intended semantics
  was direction, so long vectors were being pushed down the list. It would be
  confirmed by finding that the documents that moved are the ones whose
  embeddings had unusual lengths.}
\item Why can a dot product of zero not tell you that two documents are
  unrelated?
  \answerto{Because it is a statement about the angle between two vectors, and
  whether the angle tracks relatedness is a fact about how the vectors were
  produced. Program~\ref{prog:F08} makes the same point about high
  similarities.}
\item Give a pair of two-dimensional vectors for which the cosine similarity
  is high and the Euclidean distance is large, and say in one sentence how you
  built it.
  \answerto{Any two nearly parallel vectors of very different lengths, for
  example $(1, 0)$ and $(100, 1)$. Take a small angle to fix the cosine, then
  stretch one of them to make the distance whatever you like.}
\item A claim about $\val{f09.dim}$-dimensional embeddings is supported by a
  diagram of two arrows. Give the one question to ask about the claim.
  \answerto{Is it an identity, or is it about how the vectors are arranged? An
  identity holds at any number of components; an arrangement claim needs
  measurement in the space it is about.}
\item This program defined every operation on components first and drew the
  arrow second. Say what would have gone wrong had it been written the other
  way round.
  \answerto{Every definition would have rested on a picture available at two
  and three components, so nothing would have extended past three without
  being defined again \dash{} and the reader would have no way to tell which
  of their intuitions came from the definition and which from the drawing.}
\end{furtherproblems}
